        %auto-ignore
\section{What Causes Hyperparameter Transfer?}\label{sec:cause_hp}

In a popular misconception, hyperparameter transfer is implied by the existence of a limit.
For example, the fact that $\mu$P transfers hyperparameters, in this misconception, is because of the existence of the feature learning limit (aka the $\mu$ limit), the limit of $\mu$P as width goes to infinity.
However, this is not the case.
Indeed, there are a plethora of infinite-width limits, such as the NTK limit, but there can only be one way how the optimal hyperparameters scale, so existence cannot imply transfer.
In a stronger version of this misconception, transfer is implied by the existence of a ``feature learning'' limit. But again, this is False, because there are infinite number of feature learning limits (where the $\mu$ limit is the unique maximal one).

Instead, what is true is that the \emph{optimal} limit implies the transfer of \emph{optimal} hyperparameters.
For example, in the width limit case, $\mu$P is the unique parametrization that yields a maximal feature learning limit. Compared to all other limits, this is obviously the optimal one. Hence $\mu$P can transfer hyperparameters across width.

So far, there is no \emph{a priori} definition for the ``optimality'' of a limit: One can only tell by \emph{classifying} all possible limits; it turns out only a small number of different behavior can occur in the limit, and thus one can manually inspect for which limit is the optimal one.

Similarly, in this work, to \emph{derive} a depthwise scaling that allows transfer, we need to \emph{classify} all possible infinite depth limits --- and Depth-$\mu$P will turn out to be optimal in a sense that we define later in the paper.%
\footnote{There are important nuances here that will be spelled out in an upcoming paper. For example, if the space of hyperparameters is not chosen correctly, then it could appear that no limit is \emph{optimal} in any manner. For example, if one in (widthwise) SP, one only thinks about the 1D space of the global learning rate, then all infinite-width limits are defective --- and indeed there is no hyperparameter transfer where the bigger always does better.}
More interestingly than the width case, here we have multiple modes of feature learning when taking the depth limit and it is important to discern which mode of feature learning is optimal.
Thus, again, it is \emph{insufficient} to derive any one limit, even with feature learning, and be able to infer it yields HP transfer.

In \cref{sec:exp}, we provide experiments with $1/L$ block scaling $(\alpha, \gamma) = (1, 0)$, aka ODE scaling, which provably induces feature learning in the infinite-depth limit, but is sub-optimal. Our results show a significant shift in the optimal learning rate with this parametrization.